\section{Complete small-block proofs through five inputs}\label{app:small}

This appendix supplies the detailed algebraic constructions used before the
six-input theorem.  Every displayed product combines disjoint blocks, and
all signs refer to a choice of sum or difference.

\subsection{One and two inputs}

For one input, divisibility by $1$ is immediate.  For two arbitrary inputs
$a,b$, if they have the same parity then $a-b$ is even; otherwise $ab$ is
even.  This proves $P(2;2)$.  If one input is divisible by $3$, their product
is divisible by $3$; otherwise both are nonzero modulo $3$ and a suitable
$a\pm b$ is divisible by $3$.  Hence $P(2;3)$.

If $a,b$ are odd, one of $a+b,a-b$ is divisible by $4$.  If $a,b$ are even,
then the preceding factor-$3$ construction is automatically even, proving a
$6$-multiple.  For a multiple of $8$, if one input is divisible by $4$ use
the product; otherwise both are $2\pmod4$, and according as their residues
modulo $8$ agree or are opposite, their difference or sum is divisible by
$8$.

\subsection{Three arbitrary inputs}

\begin{theorem}\label{app:three7}
Every three integers admit an expression divisible by $7$.
\end{theorem}
\begin{proof}
If one input is divisible by $7$, multiply it by any legal expression on the
other two leaves.  Otherwise the residues lie in the three nonzero signed
classes $\{\pm1\},\{\pm2\},\{\pm3\}$ modulo $7$.  A collision gives a sum
or difference divisible by $7$; multiply that block by the remaining third
leaf.  If all three classes occur, relabel and choose signed representatives
so that $a\equiv1$, $b\equiv2$, $c\equiv3\pmod7$.  Since
$bc\equiv-1\pmod7$, one of the legal expressions $bc+a,bc-a,a-bc$ is
divisible by $7$ after restoring the selected signs.
\end{proof}

\begin{theorem}\label{app:three8}
Every three integers admit an expression divisible by $8$.
\end{theorem}
\begin{proof}
If all three are even, use the product.  If exactly two, $x,y$, are even and
$z$ is odd, use the product when one even input is divisible by $4$.
Otherwise $x,y\equiv2\pmod4$, and one of $x+y,x-y$ is divisible by $8$;
multiply it by $z$.

If exactly one input $e$ is even and $a,b$ are odd, then when $4\mid e$ use
$e(a+b)$, while when $e\equiv2\pmod4$ choose $a\pm b$ divisible by $4$ and
use $e(a\pm b)$.  Finally, for three odd inputs, the residues occupy the two
signed classes $\{\pm1\},\{\pm3\}$ modulo $8$; a colliding pair has a sum
or difference divisible by $8$, which is multiplied by the third input.
\end{proof}

\begin{theorem}\label{app:three9}
Every three integers admit an expression divisible by $9$.
\end{theorem}
\begin{proof}
If one input is divisible by $9$, use the product.  If at least two inputs are
divisible by $3$ but none by $9$, those two are in the signed class
$\{\pm3\}\pmod9$, so their suitable sum or difference is divisible by $9$;
multiply that block by the remaining third leaf.  If exactly one input $c$ is
divisible by $3$, the other two are nonzero modulo $3$ and a suitable
$a\pm b$ is divisible by $3$; hence $9\mid c(a\pm b)$.

It remains that no input is divisible by $3$.  The nonzero signed classes
modulo $9$ are
\[
 \{\pm1\},\quad\{\pm2\},\quad\{\pm4\}.
\]
A collision gives a $9$-multiple pair, which is multiplied by the remaining
third leaf.  If all three classes occur, normalize to $a\equiv1$,
$b\equiv2$, $c\equiv4\pmod9$.  Then $bc\equiv-1\pmod9$, so a suitable
$bc\pm a$ is divisible by $9$.
\end{proof}

\begin{theorem}\label{app:three10}
Every three integers admit an expression divisible by $10$.
\end{theorem}
\begin{proof}
If one input, say $a$, is divisible by $5$ and at least one of the three
inputs is even, the full product $abc$ is divisible by $10$.  If all three
inputs are odd, use $a(b+c)$, which is divisible by $10$ because $b+c$ is
even.

Assume no input is divisible by $5$.  Two inputs lie in one of the signed
classes $\{\pm1\},\{\pm2\}$ modulo $5$; call them $a,b$.  Choose a sign so
that $5\mid a\pm b$.  If $a,b$ have the same parity, this block is divisible
by $10$ and is multiplied by the third input $c$.  If they have opposite
parity and $c$ is even, then $c(a\pm b)$ is divisible by $10$.

The only remaining parity pattern has $a,b$ of opposite parity and $c$ odd.
If the two odd inputs lie in the same signed class modulo $5$, their suitable
sum or difference is divisible by $10$ and is multiplied by the remaining
even input.  Otherwise, after sign normalization, the two inputs in one signed
class contribute $r,r$ and the odd input in the other class contributes
$-2r$ modulo $5$.  Their signed three-term sum is divisible by $5$ and is
even (odd $+$ even $-$ odd), hence divisible by $10$.
\end{proof}

\begin{theorem}[Three-input $12$ theorem]\label{app:three12}
Every three integers admit an expression divisible by $12$.
\end{theorem}
\begin{proof}
If all three are even and one is divisible by $3$, their product is divisible
by $24$.  If none is divisible by $3$, choose $3\mid a\pm b$; the block
$a\pm b$ together with the remaining even factor supplies a factor $4$.

Suppose exactly two inputs $x,y$ are even and $z$ is odd.  If any input is
divisible by $3$, the product contains $4\cdot3$.  Otherwise, if one even
input, say $x$, is divisible by $4$, choose $3\mid y\pm z$ and use
$x(y\pm z)$.  If neither even input is divisible by $4$, then
$x\equiv y\equiv2\pmod4$; choose $3\mid x\pm y$.  Either sign is divisible
by $4$, so $(x\pm y)z$ is divisible by $12$.

Suppose exactly one input $e$ is even and $a,b$ are odd.  If $4\mid e$, then
either the product already contains a factor $3$, or a suitable $a\pm b$ is
divisible by $3$; in either case a $12$-multiple results.  Now let
$e\equiv2\pmod4$.  If $3\mid e$, choose $a\pm b$ divisible by $4$.  If
neither odd input is divisible by $3$, choose $a\pm b$ divisible by $3$;
this block is even and $e$ supplies the second factor $2$.  If both odd
inputs are divisible by $3$, choose $a\pm b$ divisible by $4$; it is also
divisible by $3$.

In the remaining subcase exactly one odd input, say $a$, is divisible by
$3$, while $b,e$ are nonzero modulo $3$.  Choose the sign on $b$ so that
$e\pm b\equiv0\pmod3$.  The number $e\pm b$ is odd.  Since $a$ is odd, one
of $(e\pm b)+a$ and $(e\pm b)-a$ is divisible by $4$.  Changing the sign of
$a$ does not affect divisibility modulo $3$, so the resulting signed sum is
divisible by $12$.

Finally suppose all three inputs are odd.  If one is divisible by $3$, use it
as a factor and combine the other two into a multiple of $4$.  If none is
divisible by $3$, the residues lie in the two signed classes
$\{\pm1\},\{\pm5\}$ modulo $12$; two inputs have a sum or difference
divisible by $12$, which is multiplied by the third input.
\end{proof}

\subsection{Three-input parity improvements}

\begin{theorem}[Same-parity $18$ theorem]\label{app:three18}
Any three integers of the same parity admit an expression divisible by $18$.
\end{theorem}
\begin{proof}
Work modulo $9$.  If one input $c$ is divisible by $9$, then
$c(a+b)$ is divisible by $18$ because $a+b$ is even.  If at least two inputs
are divisible by $3$ but none by $9$, those two are congruent to $\pm3$
modulo $9$; a suitable sum or difference is divisible by $9$ and, by same
parity, by $2$.  Multiply this $18$-multiple block by the remaining third
leaf.

If exactly one input $c$ is divisible by $3$, the other two are nonzero
modulo $3$, so choose $3\mid a\pm b$.  Then $9\mid c(a\pm b)$, and
$a\pm b$ is even.  If no input is divisible by $3$, use the signed classes
$\{\pm1\},\{\pm2\},\{\pm4\}$ modulo $9$.  A collision gives an even
$9$-multiple, which is multiplied by the remaining third leaf.  If all three
classes occur, normalize to $a\equiv1,b\equiv2,c\equiv4\pmod9$; then
$bc\equiv-1$, so $bc\pm a$ can be made divisible by $9$.  Since all inputs
have the same parity, $bc$ and $a$ have the same parity, so the expression is
even.
\end{proof}

\begin{lemma}\label{app:three48lemma}
For any integers $x,y,z$, either $6\mid xyz$, or after relabelling there is a
sign such that $12\mid(x\pm y)z$.
\end{lemma}
\begin{proof}
Assume $6\nmid xyz$.  If $xyz$ is odd and one input is divisible by $3$, use
it as $z$ and choose $x\pm y$ divisible by $4$.  If none is divisible by
$3$, two lie in one of the signed classes $\{\pm1\},\{\pm5\}$ modulo $12$,
and their suitable sum or difference is divisible by $12$.

If $xyz$ is even, no input is divisible by $3$.  If one input is divisible
by $4$, use it as $z$ and choose $x\pm y$ divisible by $3$.  Otherwise every
even input is $2\pmod4$.  With exactly one even input, use it as $z$ and
choose the two odd inputs so that $3\mid x\pm y$; the latter block is also
even.  With at least two even inputs, choose two as $x,y$ and make
$3\mid x\pm y$; their sum or difference is automatically divisible by $4$.
\end{proof}

\begin{theorem}[Three even inputs give $48$]\label{app:three48}
Any three even integers admit an expression divisible by $48$.
\end{theorem}
\begin{proof}
Write the inputs as $2x,2y,2z$.  If $6\mid xyz$, their product is divisible
by $48$.  Otherwise \Cref{app:three48lemma} gives
$12\mid(x\pm y)z$, and
\[
 (2x\pm2y)(2z)=4(x\pm y)z
\]
is divisible by $48$.
\end{proof}

The proof of \Cref{app:three7} also improves to $14$ when the three inputs
have one parity.  If one input is divisible by $7$, multiply it by the sum of
the other two same-parity leaves, which is even.  In the collision branch the
$7$-divisible sum or difference of the colliding same-parity pair is even and
is multiplied by the third leaf.  In the three-class branch, the product of
two same-parity inputs and the remaining input have the same parity, so the
constructed $bc\pm a$ expression is even.

\subsection{Four-input universal and parity blocks}

Besides the $24$ and $36$ theorems proved in the main text, every four-input
set yields $10$, $14$, and $16$.  The modulus $10$ follows by applying
\Cref{app:three10} to three leaves and multiplying by the fourth.  For $14$,
if an even input is present, reserve it and construct a $7$-multiple from the
other three; if all four are odd, construct a $14$-multiple from three odd
leaves and multiply by the last odd leaf.

For $16$, classify by the number of odd inputs.  Four even inputs have product
divisible by $16$.  With one odd input, the three even inputs yield $48$.
With two odd and two even inputs, the odd pair yields $4$ and the product of
the even pair yields $4$.  With three odd inputs, those three yield an
$8$-multiple and the remaining even leaf supplies another factor $2$.  With
four odd inputs, split into two pairs, each yielding $4$.

\begin{theorem}[Four-input $48$ refinements]\label{app:four48}
A $48$-multiple is guaranteed for four inputs containing exactly $0,1,2$, or
$4$ odd integers.
\end{theorem}
\begin{proof}
For four even inputs, first choose an $8$-block.  If one input is divisible
by $4$, multiply it by any other even input; their product is divisible by
$8$.  If no input is divisible by $4$, then all four inputs are congruent to
$2$ modulo $4$, hence each is congruent to $2$ or $6$ modulo $8$.  They all
lie in the single signed class $\{\pm2\}$ modulo $8$, so any selected pair
has a suitable sum or difference divisible by $8$.  The two remaining even
inputs can be combined into a multiple of $3$ by \Cref{lem:pair3}; that
expression is automatically even, hence is a $6$-block.  The two disjoint
blocks give $8\cdot6=48$.  With exactly one odd input, the three even leaves
yield $48$ and the odd leaf is multiplied afterward.

For two odd inputs $a,b$ and two even inputs $x,y$, if one even input is
divisible by $4$, reserve it and apply the three-input $12$ theorem to the
other leaves.  Otherwise $x\equiv y\equiv2\pmod4$.  If one even input is
divisible by $3$, then $xy$ is divisible by $12$, while $a\pm b$ can be
made divisible by $4$.  If neither even input is divisible by $3$, choose
$x\pm y$ divisible by $3$; the same block is divisible by $12$, and again
choose $a\pm b$ divisible by $4$.

For four odd inputs, the residues modulo $12$ occupy the signed classes
$\{\pm1\},\{\pm3\},\{\pm5\}$.  Two inputs yield a $12$-multiple, while
the remaining odd pair yields a $4$-multiple.
\end{proof}

\begin{theorem}[Four even inputs give $72$]\label{app:four72}
Any four even integers admit an expression divisible by $72$.
\end{theorem}
\begin{proof}
If one input is divisible by $3$, it supplies $6$ and the remaining three
even inputs supply $12$.  Otherwise all four inputs are nonzero modulo $3$.
Modulo $12$, even nonmultiples of $3$ occupy the signed classes
$\{\pm2\},\{\pm4\}$, so two yield a $12$-multiple.  The remaining even
pair yields a multiple of $3$, automatically even, hence a $6$-multiple.
\end{proof}

\begin{theorem}[Four odd inputs give $60$]\label{app:fourodd60}
Any four odd integers admit an expression divisible by $60$.
\end{theorem}
\begin{proof}
If one input is divisible by $5$, multiply it by a $12$-multiple from the
other three odd inputs.  Suppose no input is divisible by $5$.  If two inputs
lie in one signed class modulo $20$, their suitable sum or difference is
divisible by $20$, while the remaining pair yields a multiple of $3$.

It remains that the four inputs occupy the four signed classes
$\{\pm1\},\{\pm3\},\{\pm7\},\{\pm9\}$ modulo $20$.  Normalize and label
them $a,b,c,d$ in that order.  The pair $b,c$ has a suitable sum or
difference divisible by $10$.  If $a,d$ are both divisible by $3$ or both
nonzero modulo $3$, they yield a $6$-multiple, completing $6\cdot10$.

If exactly one of $a,d$ is divisible by $3$, suppose first $3\mid a$.
Modulo $20$, the product of the classes of $b$ and $d$ is the class of $c$,
so a suitable $bd\pm c$ is divisible by $20$; multiplication by $a$ gives
$60$.  If $3\mid d$, similarly the class of $bc$ equals that of $a$, so a
suitable $bc\pm a$ is divisible by $20$ and multiplication by $d$ completes
the proof.
\end{proof}

\begin{theorem}[One odd and three even inputs give $40$]\label{app:oneodd40}
Any one odd and three even inputs admit an expression divisible by $40$.
\end{theorem}
\begin{proof}
Let the odd input be $a$ and the even inputs $x,y,z$.  If $5\mid a$, multiply
$a$ by an $8$-multiple from $x,y,z$.  If an even input is divisible by $5$,
it supplies $10$ and the other three inputs supply $8$.

Assume no input is divisible by $5$.  If two even inputs are in the same
signed class modulo $5$ and have the same residue modulo $4$, their suitable
sum or difference is divisible by $20$; the remaining odd-even product is
even.

Otherwise the three even inputs have signed-class distribution $2+1$ modulo
$5$ (the $3+0$ distribution would contain a pair with the same residue
modulo $4$).  Let $x,y$ be the majority-class pair and $z$ the singleton.
Then $x,y$ have different residues modulo $4$.  If $4\mid z$, the pair
$x,y$ yields a $10$-multiple and $az$ yields a $4$-multiple.  Finally let
$z\equiv2\pmod4$.  In the majority pair one input is $0\pmod4$ and the other
$2\pmod4$.  Pair the odd input $a$ with the $2\pmod4$ even input in the same
modulo-$5$ signed class.  If $a$ lies in the majority class, use the
$2\pmod4$ member of the majority pair.  If $a$ does not lie in the majority
class, then, because the nonzero residues modulo $5$ form only two signed
classes, $a$ must lie in the singleton class; hence $a$ and $z$ lie in the
same class.  Their suitable sum or difference is divisible by $5$.
The two unused even inputs have product divisible by $8$.  Thus the product
is divisible by $40$.
\end{proof}

\begin{theorem}[Two odd and two even inputs give $40$]\label{app:twoodd40}
Any two odd and two even inputs admit an expression divisible by $40$.
\end{theorem}
\begin{proof}
Let the odd inputs be $p,q$ and the even inputs $x,y$.  If one even input is
divisible by $4$, multiply it by a $10$-multiple from the remaining three
leaves.  Hence assume $x\equiv y\equiv2\pmod4$.

If some input is divisible by $5$, reserve its factor $5$ (or $10$ for an
even input) and use the three-input $8$ theorem on the other leaves.  Now
assume all four inputs are nonzero modulo $5$.

If $p,q$ occupy one signed class modulo $5$, their suitable sum or difference
is divisible by $10$, and $xy$ is divisible by $4$.  If $x,y$ occupy one
signed class modulo $5$, their suitable sum or difference is divisible by
$5$ and by $4$, while the odd pair yields $4$.

It remains that each parity pair occupies the two different signed classes
modulo $5$.  Relabel so that $p,x$ occupy one class and $q,y$ the other.
Consider
\[
 E_1=px\pm qy,\qquad E_2=py\pm qx.
\]
Choose signed representatives modulo $5$ so that
\[
 p\equiv\varepsilon_p r,\qquad x\equiv\varepsilon_x r,\qquad
 q\equiv2\varepsilon_q r,\qquad y\equiv2\varepsilon_y r,
\]
where each $\varepsilon$ is $+1$ or $-1$.  Put
$\delta=\varepsilon_p\varepsilon_q\varepsilon_x\varepsilon_y$.  Direct
multiplication of these congruences shows that $E_1$ is divisible by $5$
when its sign is $\delta$, whereas $E_2$ is divisible by $5$ when its sign
is $-\delta$.  Thus the modulo-$5$ signs required by $E_1$ and $E_2$ are
opposite.

Now write $x=2X$ and $y=2Y$, where $X,Y$ are odd.  For a sign
$s\in\{\pm1\}$,
\[
 E_1=2(pX+s qY),\qquad E_2=2(pY+s qX).
\]
Divisibility by $8$ is therefore equivalent to divisibility by $4$ of the
corresponding parenthesized odd sum or difference.  The two pairs of odd
terms satisfy
\[
 (pX)(qY)=(pY)(qX)=pqXY.
\]
For two odd integers modulo $4$, the sign that makes their sum or difference
zero modulo $4$ depends only on whether their residues are equal or opposite.
The displayed product identity shows that the pair $pX,qY$ is equal or
opposite modulo $4$ in exactly the same way as the pair $pY,qX$.  Hence
$E_1$ and $E_2$ require the same sign for divisibility by $8$.

The modulo-$5$ signs for $E_1$ and $E_2$ are opposite, whereas their
modulo-$8$ signs are the same.  Therefore, for exactly one of $E_1,E_2$, the
same choice of sign gives divisibility by both $5$ and $8$.  That legal
read-once expression is divisible by $40$.
\end{proof}

The theorem $3\OO1\EE\to60$ is proved separately in
\Cref{app:mixed60}.  The simpler $3\OO1\EE\to20$ follows by constructing a
$10$-multiple from the three odd inputs and multiplying by the even input.

\subsection{Five-input parity theorems}

\begin{theorem}[Five odd inputs]\label{app:fiveodd}
Five odd inputs admit expressions divisible by $180$ and by $240$.
\end{theorem}
\begin{proof}
Suppose first that no input is divisible by $5$.  The five inputs occupy the
four nonzero odd signed classes modulo $20$, so a pair yields a $20$-multiple.
The remaining three odd inputs yield an $18$-multiple, giving $360$, hence
$180$; using instead their universal $12$-multiple gives $240$.

If one input $a$ is divisible by $5$, among the other four odd inputs choose
a pair $b,c$ whose sum or difference is divisible by $12$.  Then
$a(b\pm c)$ is divisible by $60$.  The remaining odd pair can be combined
into a multiple of $3$, giving $180$, or into a multiple of $4$, giving
$240$.
\end{proof}

\begin{theorem}[Five even inputs]\label{app:fiveeven}
Five even inputs admit an expression divisible by $480$.
\end{theorem}
\begin{proof}
If one input is divisible by $5$, it is divisible by $10$; the remaining
four even inputs yield $48$.  Otherwise the five inputs occupy the two
nonzero signed classes modulo $5$, so two have a suitable sum or difference
divisible by $5$ and by $2$, hence by $10$.  The remaining three even inputs
yield $48$.  In both branches the disjoint blocks give $480$.
\end{proof}
